\section{Introduction}
\label{sec:intro}

Score-based diffusion models generate data by reversing a noising process~\cite{song2021sde,ho2020ddpm}.
Critically damped Langevin diffusion (CLD)~\cite{dockhorn2022cld} runs this process in phase space: each data coordinate $x$ is paired with a velocity $v$, the pair is driven by a damped oscillator flow, and noise enters through the velocity only.
\emph{The standard CLD prescribes its noise through its stochastic differential equation (SDE). It therefore never meets the threshold studied in this paper.}
Our subject is a different design choice: keeping the CLD flow but specifying the noise schedule by its marginals, as an isotropic blur of width $\sigma(t)=\smin e^{rt}$ that grows at a fixed rate $r$.

Why consider such schedules?
Specifying a schedule by its marginals, as in variance-exploding and EDM-type diffusion models~\cite{song2021sde,karras2022edm}, fixes the noise level $\sigma$ directly; the preconditioning of the network inputs and the choice of time steps can then be taken over from those designs.
Making the blur isotropic puts positions and velocities on the same scale, so a single $\sigma$ states how much phase-space detail has been blurred away.
We do not test here whether these conveniences pay off in practice for CLD.
This paper only establishes what happens to the reverse samplers when such a schedule widens more slowly than the flow stretches.

For such schedules the noise required to produce the marginals is $G=2\sigma^2(rI-\Sym B)$, where $B$ is the flow matrix.
It is positive semidefinite only if $r\geq\kappa=\lambda_{\max}(\Sym B)$, the fastest stretching rate of the flow~\cite{namba2026threshold}.
Below $\kappa$, the standard reverse SDE~\cite{anderson1982reverse,haussmann1986time} does not exist.
Its natural repair, which clips the negative part of $G$, is biased, whereas a corrected sampler and the probability-flow ODE~\cite{song2021sde} stay exact.
These statements are proved in Ref.~\cite{namba2026threshold}; the threshold itself is the phase-space form of a criterion for the coarse-grained arrow of time~\cite{namba2026arrow}.

This paper asks the practical questions that the proofs leave open, for the actual CLD flow:
\emph{where} the clipped sampler breaks, \emph{how much}, and whether the corrected sampler still fixes it when the score is \emph{learned}.
Our answers are as follows.
The measured threshold sits where the theory puts it (Sec.~\ref{sec:exact}).
The damage is made at the end of sampling, while the blur is close to $\smin$, and it is large only in specific regions of parameter space (Sec.~\ref{sec:map}).
With learned scores, the clipped sampler breaks and the corrected sampler fixes it, the score error simply adding to the threshold effect (Sec.~\ref{sec:learned}).
The probability-flow ODE carries score errors that the corrected SDE forgets (Sec.~\ref{sec:forgetting}).
In typical settings the effect is a few per cent (Sec.~\ref{sec:discussion}).

Throughout, we mark the status of every statement:
a \emph{Proposition} is proved;
a \emph{Derived result} is an exact or asymptotic formula derived here and checked against direct numerical solution;
an \emph{Observation} is measured in simulations, with seeds and error bars;
a \emph{Conjecture} has not been tested.
